\documentclass[11pt]{article}
\usepackage[T1]{fontenc}
\usepackage[utf8]{inputenc}
\usepackage{lmodern}
\usepackage{geometry}
\usepackage{hyperref}
\usepackage{amsmath}
\usepackage{amssymb}
\usepackage{enumitem}
\geometry{margin=1in}

\title{Block-Diagonal Jacobian of the LDG Residual in Exasim}
\author{}
\date{}

\begin{document}
\maketitle

\section*{1. Summary}

For the LDG path in Exasim, the exact block-diagonal Jacobian of

\[
J(U) = \frac{\partial R(U)}{\partial U}
\]

with $U$ the global vector of the degrees of freedom of $u_h$, is most
naturally implemented as one dense element block per element:

\[
J_{\mathrm{diag}} = \operatorname{diag}(J_{11}, J_{22}, \dots, J_{N_e N_e}),
\]

where each block

\[
J_{KK} = \frac{\partial R_K}{\partial U_K}
\]

has size $(npe \cdot ncu) \times (npe \cdot ncu)$, with $R_K$ the residual
rows belonging to element $K$.

In the current Exasim backend, the cleanest exact implementation is not to
derive and assemble a separate sparse Jacobian by hand, but to reuse the
existing residual tangent path:

\begin{itemize}[leftmargin=*]
\item \texttt{backend/Discretization/residual.cpp}
\item \texttt{backend/Discretization/qresidual.cpp}
\item \texttt{backend/Discretization/uresidual.cpp}
\item \texttt{backend/Discretization/getuhat.cpp}
\item \texttt{backend/Discretization/matvec.cpp}
\end{itemize}

Specifically:

\begin{itemize}[leftmargin=*]
\item if Enzyme AD is enabled, assemble each diagonal block column-by-column
using \texttt{dResidual(...)}
\item otherwise, assemble each block column-by-column using
\texttt{MatVec(...)} with finite differences
\end{itemize}

That approach automatically includes all LDG chain-rule contributions through:

\begin{itemize}[leftmargin=*]
\item \texttt{GetUhat(...)}
\item \texttt{GetQ(...)}
\item \texttt{GetW(...)}
\item \texttt{GetAv(...)}
\item \texttt{RuElem(...)}
\item \texttt{RuFace(...)}
\end{itemize}

\section*{2. Assumptions}

\begin{itemize}[leftmargin=*]
\item $U$ contains only the DOFs of $u_h$, not the auxiliary variables $q$,
$w$, or face trace \texttt{uh}.
\item The residual $R(U)$ is exactly the one returned by
\texttt{Residual(sol, res, app, driver\_abi, master, mesh, tmp, common, handle, backend)}
in \texttt{backend/Discretization/residual.cpp}.
\item Therefore $q$, \texttt{uh}, \texttt{wdg}, and optional AV-related fields
are implicit functions of $U$, because \texttt{Residual(...)} recomputes them
internally before assembling \texttt{res.Ru}.
\item A ``block'' means one element block of size:
\[
n_{\mathrm{loc}} = npe \cdot ncu.
\]
\item The desired matrix is the exact block-diagonal of the full LDG Jacobian,
not a simplified ``volume-only'' or ``element-local-only'' approximation.
\item I am focusing on the LDG residual path, not the HDG Schur-complement
system.
\item I assume the block-diagonal Jacobian may later be inverted batched, for
example for a block-Jacobi preconditioner.
\end{itemize}

\section*{3. Design reasoning}

The key point is that the residual returned by \texttt{Residual(...)} is not
just a function of the element-local $u$. It is a composite map:

\[
U \;\longmapsto\; \hat U(U), \; Q(U), \; W(U), \; AV(U) \;\longmapsto\; R(U).
\]

So the exact diagonal block must include all chain-rule paths created by the
residual pipeline.

\subsection*{A. Residual structure used by \texttt{Residual(...)}}

In serial, the LDG residual pipeline is:

\begin{enumerate}[leftmargin=*]
\item \texttt{GetUhat(...)}
\item \texttt{GetQ(...)}
\item \texttt{GetW(...)} if \texttt{common.ncw > 0}
\item \texttt{GetAv(...)} if AV is active and unfrozen
\item \texttt{RuElem(...)}
\item \texttt{RuFace(...)}
\item \texttt{PutFaceNodes(...)}
\item final scaling by \texttt{-1} and optionally \texttt{1/common.dtfactor}
\end{enumerate}

This is in \texttt{backend/Discretization/residual.cpp}.

So if I define the raw assembled residual before the final sign/time scaling by
$\widetilde R(U)$, then the residual returned by \texttt{Residual(...)} is

\[
R(U) = \alpha \, \widetilde R(U),
\]

with

\[
\alpha =
\begin{cases}
-1, & \text{steady case},\\[4pt]
-\dfrac{1}{\texttt{common.dtfactor}}, & \text{time-dependent case}.
\end{cases}
\]

Therefore each diagonal block is

\[
J_{KK} = \alpha \, \frac{\partial \widetilde R_K}{\partial U_K}.
\]

That final scalar factor must be included.

\subsection*{B. Element-block definition}

Let $U_K \in \mathbb{R}^{n_{\mathrm{loc}}}$ be the $u_h$ DOFs of element $K$,
with

\[
n_{\mathrm{loc}} = npe \cdot ncu.
\]

Let $R_K(U)$ be the residual entries in \texttt{res.Ru} associated with
element $K$. Then

\[
J_{KK} = \frac{\partial R_K}{\partial U_K}.
\]

This is the exact diagonal block. It keeps all dependence of the residual rows
of element $K$ on the DOFs of the same element $K$, while discarding all
derivatives with respect to DOFs from other elements.

\subsection*{C. Chain rule through the LDG pipeline}

The raw residual for element $K$ can be viewed schematically as

\[
\widetilde R_K(U)
=
\widetilde R^{\mathrm{elem}}_K\big(U_K, Q_K(U), W_K(U), AV_K(U)\big)
+
\sum_{f \subset \partial K}
\widetilde R^{\mathrm{face}}_{Kf}
\big(U, Q(U), W(U), \hat U(U)\big).
\]

Hence the exact diagonal block is

\[
\frac{\partial \widetilde R_K}{\partial U_K}
=
\frac{\partial \widetilde R^{\mathrm{elem}}_K}{\partial U_K}
+
\frac{\partial \widetilde R^{\mathrm{elem}}_K}{\partial Q_K}\frac{\partial Q_K}{\partial U_K}
+
\frac{\partial \widetilde R^{\mathrm{elem}}_K}{\partial W_K}\frac{\partial W_K}{\partial U_K}
+
\frac{\partial \widetilde R^{\mathrm{elem}}_K}{\partial AV_K}\frac{\partial AV_K}{\partial U_K}
+
\sum_{f \subset \partial K}
\frac{\partial \widetilde R^{\mathrm{face}}_{Kf}}{\partial U_K}^{\mathrm{total}}.
\]

The face contribution expands further into:

\[
\frac{\partial \widetilde R^{\mathrm{face}}_{Kf}}{\partial U_K}^{\mathrm{total}}
=
\frac{\partial \widetilde R^{\mathrm{face}}_{Kf}}{\partial U_{K,f}}
+
\frac{\partial \widetilde R^{\mathrm{face}}_{Kf}}{\partial Q_{K,f}}
\frac{\partial Q_{K,f}}{\partial U_K}
+
\frac{\partial \widetilde R^{\mathrm{face}}_{Kf}}{\partial W_{K,f}}
\frac{\partial W_{K,f}}{\partial U_K}
+
\frac{\partial \widetilde R^{\mathrm{face}}_{Kf}}{\partial \hat U_f}
\frac{\partial \hat U_f}{\partial U_K}.
\]

This is the correct mathematical content of the diagonal block in Exasim.

\subsection*{D. Contribution from \texttt{GetUhat(...)}}

\texttt{GetUhat(...)} is in \texttt{backend/Discretization/getuhat.cpp}.

For LDG:

\begin{itemize}[leftmargin=*]
\item on interior faces, \texttt{sol.uh} is copied from one side's $u$
\item on boundary faces, \texttt{sol.uh} is produced by \texttt{UbouDriver(...)}
\end{itemize}

Therefore the derivative

\[
\frac{\partial \hat U_f}{\partial U_K}
\]

can be nonzero for:

\begin{itemize}[leftmargin=*]
\item the element that owns the copied interior-face trace
\item any boundary face of element $K$
\end{itemize}

This derivative is implemented in the tangent path by:

\begin{itemize}[leftmargin=*]
\item \texttt{GetdUhat(...)}
\item \texttt{dUhatBlock(...)}
\end{itemize}

in the same file.

That means the exact diagonal block must include the way a perturbation in
$U_K$ changes \texttt{uh} first, and then changes face residuals and
$q$-recovery through that \texttt{uh}.

\subsection*{E. Contribution from \texttt{GetQ(...)}}

\texttt{GetQ(...)} is in \texttt{backend/Discretization/residual.cpp}.

It computes $q$ by:

\begin{itemize}[leftmargin=*]
\item \texttt{RqElem(...)}
\item \texttt{RqFace(...)}
\item \texttt{PutFaceNodes(...)}
\item \texttt{ApplyMinv(...)}
\item \texttt{ArrayInsert(...)} into \texttt{sol.udg}
\end{itemize}

From the implementation in \texttt{backend/Discretization/qresidual.cpp}, the
elementwise operator is

\[
Q_K = M_K^{-1}\left(C_K U_K - E_K \hat U_K + S^q_K\right),
\]

where:

\begin{itemize}[leftmargin=*]
\item $M_K$ is the local mass matrix inverse applied through
\texttt{ApplyMinv(...)}
\item $C_K$ is the element gradient operator from \texttt{RqElemBlock(...)}
\item $E_K$ is the face lifting operator from \texttt{RqFaceBlock(...)}
\item $S^q_K$ is the wave-only source contribution when
\texttt{common.wave == 1}
\end{itemize}

Therefore the exact local derivative is

\[
\frac{\partial Q_K}{\partial U_K}
=
M_K^{-1}\left(
C_K - E_K \frac{\partial \hat U_K}{\partial U_K}
\right),
\]

plus any wave-specific term if active.

In the current backend, this derivative is not assembled explicitly as a
matrix in the LDG path. Instead it is applied as a tangent operator by:

\begin{itemize}[leftmargin=*]
\item \texttt{GetdQ(...)}
\item \texttt{dRqElem(...)}
\item \texttt{dRqFace(...)}
\item \texttt{ApplyMinv(...)}
\end{itemize}

in \texttt{backend/Discretization/residual.cpp}.

So \texttt{GetdQ(...)} is exactly the chain-rule implementation of
$dQ/dU$.

\subsection*{F. Contribution from \texttt{GetW(...)}}

If \texttt{common.ncw > 0}, \texttt{GetW(...)} in
\texttt{backend/Discretization/residual.cpp} computes the auxiliary field
\texttt{wdg}.

This can happen through:

\begin{itemize}[leftmargin=*]
\item explicit time/update formulas
\item Newton solve of \texttt{Eos(w,u)=0}
\item sourcew/DAE updates
\end{itemize}

Therefore the residual also contains a chain rule term

\[
\frac{\partial \widetilde R_K}{\partial W_K}\frac{\partial W_K}{\partial U_K}.
\]

This term is conceptually required in the exact block diagonal.

However, the current AD tangent path in \texttt{dRuResidual(...)} contains:

\begin{itemize}[leftmargin=*]
\item \texttt{// TODO: DAE functionality not implemented yet}
\end{itemize}

So for cases with nontrivial $w$ evolution, the exactness of the AD-based
block assembly depends on the specific \texttt{GetW(...)} branch. This is an
important implementation limitation.

\subsection*{G. Contribution from \texttt{RuElem(...)}}

\texttt{RuElemBlock(...)} in \texttt{backend/Discretization/uresidual.cpp}
computes the element-volume contribution.

At Gauss points it evaluates:

\begin{itemize}[leftmargin=*]
\item time contribution
\item \texttt{SourceDriver(...)}
\item \texttt{FluxDriver(...)}
\end{itemize}

then applies geometry and basis testing through:

\begin{itemize}[leftmargin=*]
\item \texttt{ApplyXx4(...)}
\item \texttt{Gauss2Node(...)}
\end{itemize}

So the elemental Jacobian contribution is

\[
\frac{\partial \widetilde R_K^{\mathrm{elem}}}{\partial U_K}
=
\frac{\partial \widetilde R_K^{\mathrm{elem}}}{\partial U_K}\Big|_{\text{direct}}
+
\frac{\partial \widetilde R_K^{\mathrm{elem}}}{\partial Q_K}\frac{\partial Q_K}{\partial U_K}
+
\frac{\partial \widetilde R_K^{\mathrm{elem}}}{\partial W_K}\frac{\partial W_K}{\partial U_K}
+
\frac{\partial \widetilde R_K^{\mathrm{elem}}}{\partial AV_K}\frac{\partial AV_K}{\partial U_K}.
\]

In the current backend, this total directional derivative is implemented by:

\begin{itemize}[leftmargin=*]
\item \texttt{dRuElemBlock(...)}
\end{itemize}

which consumes:

\begin{itemize}[leftmargin=*]
\item \texttt{sol.dudg}, already containing \texttt{dq} inserted by
\texttt{GetdQ(...)}
\item \texttt{sol.dwdg} if available
\item \texttt{sol.dodg} if AV is active
\end{itemize}

So \texttt{dRuElemBlock(...)} is already the exact elemental Jacobian action
needed for each diagonal block column.

\subsection*{H. Contribution from \texttt{RuFace(...)}}

\texttt{RuFaceBlock(...)} in \texttt{backend/Discretization/uresidual.cpp}
computes the face contribution by:

\begin{itemize}[leftmargin=*]
\item gathering \texttt{uh}, \texttt{udg}, \texttt{wdg} on faces
\item interpolating to face Gauss points
\item evaluating \texttt{FhatDriver(...)} on interior faces
\item evaluating \texttt{FbouDriver(...)} on boundary faces
\item applying \texttt{ApplyJacFhat(...)}
\item projecting back by \texttt{Gauss2Node(...)}
\end{itemize}

Therefore the face Jacobian contribution for the diagonal block is

\[
\sum_{f \subset \partial K}
\left(
\frac{\partial \widetilde R_{Kf}}{\partial U_{K,f}}
+
\frac{\partial \widetilde R_{Kf}}{\partial Q_{K,f}}\frac{\partial Q_{K,f}}{\partial U_K}
+
\frac{\partial \widetilde R_{Kf}}{\partial W_{K,f}}\frac{\partial W_{K,f}}{\partial U_K}
+
\frac{\partial \widetilde R_{Kf}}{\partial \hat U_f}\frac{\partial \hat U_f}{\partial U_K}
\right).
\]

This total derivative is implemented by:

\begin{itemize}[leftmargin=*]
\item \texttt{dRuFaceBlock(...)}
\end{itemize}

which uses:

\begin{itemize}[leftmargin=*]
\item \texttt{sol.duh}
\item \texttt{sol.dudg}
\item \texttt{sol.dwdg}
\end{itemize}

So again, the exact diagonal-block action is already available through the
tangent residual path.

\subsection*{I. Why the exact diagonal block should be assembled from \texttt{dResidual(...)}}

This is the most important implementation decision.

Because LDG in Exasim recomputes:

\begin{itemize}[leftmargin=*]
\item \texttt{uh}
\item \texttt{q}
\item \texttt{w}
\item AV
\end{itemize}

inside \texttt{Residual(...)}, the exact diagonal block is not just the local
derivative of \texttt{RuElemBlock(...)} with respect to the local nodal $u$.

It must include all implicit dependencies created by the preprocessing steps
inside \texttt{Residual(...)}.

Therefore, if the goal is the exact block diagonal of the Jacobian of the
actual residual returned by \texttt{Residual(...)}, the safest implementation
is:

\begin{itemize}[leftmargin=*]
\item seed a local basis vector in $u$
\item run \texttt{dResidual(...)} or \texttt{MatVec(...)}
\item extract only the same-element residual rows
\end{itemize}

That is exact by construction for the implemented residual.

\section*{4. Proposed implementation}

\subsection*{A. Affected files}

The implementation should be built around these files:

\begin{itemize}[leftmargin=*]
\item \texttt{backend/Discretization/residual.cpp}
\item \texttt{backend/Discretization/qresidual.cpp}
\item \texttt{backend/Discretization/uresidual.cpp}
\item \texttt{backend/Discretization/getuhat.cpp}
\item \texttt{backend/Discretization/matvec.cpp}
\end{itemize}

\subsection*{B. Data structure for the block-diagonal matrix}

Let

\[
n_{\mathrm{loc}} = npe \cdot ncu.
\]

Store the diagonal blocks as a dense batched array:

\[
BJ \in \mathbb{R}^{n_{\mathrm{loc}} \times n_{\mathrm{loc}} \times ne1},
\]

where \texttt{ne1} is the number of owned elements whose residual rows are in
\texttt{res.Ru}.

In code terms this should be a flat array of length:

\begin{itemize}[leftmargin=*]
\item \texttt{nloc * nloc * common.ne1}
\end{itemize}

stored element-block by element-block, column-major within each block so it is
immediately compatible with batched BLAS/LAPACK-style inversion.

\subsection*{C. Exact assembly algorithm using \texttt{dResidual(...)}}

This is the recommended exact implementation when Enzyme AD is available.

For each owned element $K$:

\begin{enumerate}[leftmargin=*]
\item Define local block size
\[
n_{\mathrm{loc}} = npe \cdot ncu.
\]

\item For each local column $j = 0, \dots, n_{\mathrm{loc}}-1$:
\begin{itemize}[leftmargin=*]
\item zero all tangent fields:
\begin{itemize}[leftmargin=*]
\item \texttt{sol.dudg}
\item \texttt{sol.dwdg} if used
\item \texttt{sol.duh}
\item \texttt{sol.dodg} if AV is active
\item \texttt{res.dRu}
\item \texttt{res.dRq}
\item \texttt{res.dRh}
\end{itemize}
\item place a unit seed in \texttt{sol.dudg} corresponding to the $j$-th
\texttt{u} DOF of element $K$
\item call:
\begin{itemize}[leftmargin=*]
\item \texttt{dResidual(sol, res, app, driver\_abi, master, mesh, tmp, common, handle, backend)}
\end{itemize}
\item extract the element-residual slice of \texttt{res.dRu} corresponding to
element $K$
\item store that slice as column $j$ of the dense block $J_{KK}$
\end{itemize}
\end{enumerate}

This works because \texttt{dResidual(...)} already performs:

\begin{itemize}[leftmargin=*]
\item \texttt{GetdUhat(...)}
\item \texttt{GetdQ(...)}
\item \texttt{GetdAv(...)}
\item \texttt{dRuElem(...)}
\item \texttt{dRuFace(...)}
\end{itemize}

and applies the final residual scaling.

So each call returns exactly:

\[
J(U)\,e^{(K,j)},
\]

and extracting the rows of element $K$ gives precisely the column $j$ of
$J_{KK}$.

\subsection*{D. Exact assembly algorithm using \texttt{MatVec(...)}}

If Enzyme AD is not available, use the same column-by-column procedure but
replace \texttt{dResidual(...)} with:

\begin{itemize}[leftmargin=*]
\item \texttt{MatVec(...)}
\end{itemize}

from \texttt{backend/Discretization/matvec.cpp}.

For each seed vector $e^{(K,j)}$:

\begin{itemize}[leftmargin=*]
\item \texttt{MatVec(...)} returns $J(U)e^{(K,j)}$ using either:
\begin{itemize}[leftmargin=*]
\item first-order finite difference
\item second-order finite difference
\end{itemize}
\item then extract the same-element residual rows and store the column
\end{itemize}

This is more expensive and approximate, but it differentiates the actual
\texttt{Residual(...)} pipeline and therefore includes all implemented
couplings.

\subsection*{E. Practical column-seeding detail}

The seed must be inserted only into the $u$ components of \texttt{sol.dudg},
not into $q$ or other fields.

That means:

\begin{itemize}[leftmargin=*]
\item within element $K$
\item within the \texttt{udg} layout
\item only the first \texttt{ncu} components per node are seeded
\end{itemize}

All $q$-component tangent entries must start from zero and be generated by
\texttt{GetdQ(...)}.

This is essential. Otherwise the result is not $\partial R/\partial U$, but a
derivative with independently perturbed $q$.

\subsection*{F. Optional inversion for block-Jacobi use}

If the goal is a block-Jacobi preconditioner, after assembling the blocks one
can invert them batched:

\[
J_{KK}^{-1}, \qquad K = 1,\dots,ne1.
\]

This can reuse the same batched inversion pattern already used elsewhere in
Exasim for dense small blocks.

The natural storage layout is chosen specifically so that a single batched
inversion call can process all element blocks.

\subsection*{G. Why this is better than hand-assembling a separate LDG Jacobian first}

A manual Jacobian assembly would need to reproduce every chain-rule dependency
already encoded in \texttt{Residual(...)}, including:

\begin{itemize}[leftmargin=*]
\item interior/boundary \texttt{uh} construction
\item local $q$-recovery
\item local $w$-recovery
\item AV differentiation
\item face orientation details
\item sign/time scaling
\end{itemize}

That is possible, but it is much more invasive and easier to get wrong.

Using \texttt{dResidual(...)} or \texttt{MatVec(...)} avoids duplicating solver
logic and guarantees consistency with the actual residual definition.

\section*{5. Risks or numerical concerns}

\begin{itemize}[leftmargin=*]
\item The AD path is not fully general for all $w$/DAE branches. In
\texttt{backend/Discretization/residual.cpp}, \texttt{dRuResidual(...)}
explicitly notes:
\begin{itemize}[leftmargin=*]
\item \texttt{TODO: DAE functionality not implemented yet}
\end{itemize}
\item Therefore:
\begin{itemize}[leftmargin=*]
\item AD-based block assembly is exact only for the branches actually covered
by the tangent implementation
\item FD-based \texttt{MatVec(...)} is more complete, but approximate
\end{itemize}
\item Finite-difference block assembly is expensive:
\begin{itemize}[leftmargin=*]
\item $n_{\mathrm{loc}}$ residual evaluations per element block
\item or one global \texttt{MatVec(...)} per local column if done naively
\end{itemize}
\item Finite-difference quality depends on:
\begin{itemize}[leftmargin=*]
\item \texttt{common.matvecTol}
\item \texttt{common.matvecOrder}
\end{itemize}
\item For interface elements in MPI, the exact diagonal block still requires
the MPI tangent/residual path, because face terms depend on exchanged neighbor
states even when only the diagonal rows are kept.
\item If AV is active and unfrozen, the exact block diagonal of the residual
must include \texttt{GetAv(...)} differentiation. Omitting that gives a
different Jacobian.
\item A hand-derived ``element-only'' approximation will generally miss
orientation-dependent face-trace effects from \texttt{GetUhat(...)}.
\end{itemize}

\section*{6. Verification plan}

\begin{enumerate}[leftmargin=*]
\item \textbf{Column test against \texttt{dResidual(...)}}
\begin{itemize}[leftmargin=*]
\item Pick an element $K$ and a random local vector $z_K$.
\item Form a global seed $z$ supported only on element $K$.
\item Compute $y = Jz$ with \texttt{dResidual(...)}.
\item Extract the element-$K$ rows of $y$.
\item Compare with $J_{KK} z_K$ from the assembled dense block.
\end{itemize}

\item \textbf{Finite-difference spot check}
\begin{itemize}[leftmargin=*]
\item For a few local columns $j$, compare the assembled column against:
\[
\frac{R(U + \varepsilon e^{(K,j)}) - R(U)}{\varepsilon}
\]
restricted to the element-$K$ residual rows.
\end{itemize}

\item \textbf{Boundary and interior element checks}
\begin{itemize}[leftmargin=*]
\item Test one boundary element and one interior element.
\item This verifies both:
\begin{itemize}[leftmargin=*]
\item \texttt{FbouDriver(...)}
\item \texttt{FhatDriver(...)}
\item \texttt{UbouDriver(...)}
\end{itemize}
\end{itemize}

\item \textbf{Physics toggles}
\begin{itemize}[leftmargin=*]
\item Verify separately for:
\begin{itemize}[leftmargin=*]
\item \texttt{common.tdep = 0}
\item \texttt{common.tdep = 1}
\item \texttt{common.ncq = 0} and \texttt{> 0}
\item \texttt{common.ncw = 0} and a case with \texttt{ncw > 0}
\item AV frozen and unfrozen
\end{itemize}
\end{itemize}

\item \textbf{Residual-consistency check}
\begin{itemize}[leftmargin=*]
\item Confirm the assembled block includes the final residual scaling by
\texttt{-1} and \texttt{1/common.dtfactor} exactly as in \texttt{Residual(...)}.
\end{itemize}

\item \textbf{Performance sanity}
\begin{itemize}[leftmargin=*]
\item Measure cost per element block.
\item If explicit block assembly is too expensive, keep the same mathematical
definition but consider:
\begin{itemize}[leftmargin=*]
\item assembling blocks only when the nonlinear state changes significantly
\item or using the block assembly only for preconditioner setup, not every
Krylov iteration
\end{itemize}
\end{itemize}
\end{enumerate}

\end{document}
